\documentclass[10pt,a4paper]{amsart}
\usepackage[margin=19mm]{geometry}
\usepackage{amsmath,amssymb,mathtools}
\usepackage[scr=rsfs,cal=euler]{mathalfa}
\usepackage{xfrac}
\usepackage[unicode,hidelinks,bookmarks=false]{hyperref}
\newcommand{\ie}{i.e.\ }
\newcommand{\cf}{cf.\ }
\newcommand{\resp}{respectively\ }
\newcommand{\Subsection}{\subsection}
\providecommand{\nobreakdash}{\nobreak\hbox{-}}
\setlength{\emergencystretch}{2em}
\multlinegap0pt
\usepackage{mdframed}
\usepackage{mdframed}
\usepackage{mdframed}
\usepackage{mdframed}
\usepackage{mathtools}
\usepackage{amssymb}
\usepackage{xcolor}
\usepackage{algorithm}\usepackage{algpseudocode}
\usepackage{mdframed}
\usepackage{float}
\usepackage[shortlabels]{enumitem}
\usepackage{cancel}
\newtheorem{assumption}{Assumption}
\newtheorem{proposition}{Proposition}
\newtheorem{remark}{Remark}
\newtheorem{lemma}{Lemma}
\newtheorem{theorem}{Theorem}
\newtheorem{corollary}{Corollary}
\newcommand{\Softmax}{\mathrm{Softmax}}
\newcommand{\Tr}{\mathrm{trace}}
\newcommand{\dee}{\mathrm{d}}
\newcommand{\mb}{\mathbb}
\newcommand{\mc}{\mathcal}
\newcommand{\tin}{t_{\mathrm{in}}}
\title{Diffusion Model's Generalization Can Be Characterized by \\ Inductive Biases toward a Data-Dependent Ridge Manifold}
\author{Ye He, Yitong Qiu, Molei Tao}
\date{Source: arXiv:2602.06021v2 (CC BY 4.0)}
\begin{document}
\maketitle

\noindent\emph{Source and licence.} Diffusion Model's Generalization Can Be Characterized by \\ Inductive Biases toward a Data-Dependent Ridge Manifold. Version arXiv:2602.06021v2 (CC BY 4.0), distributed by arXiv under the Creative Commons Attribution 4.0 International licence, http://creativecommons.org/licenses/by/4.0/, as recorded in the arXiv licence metadata for this exact version and shown on the abstract page https://arxiv.org/abs/2602.06021v2. Full licence notice: \texttt{licenses/arxiv-2602.06021-CC-BY-4.0.txt}.\par\smallskip

\noindent\textit{Editorial scope.} Counted unit: the article's theorem contraciton (source0.tex lines 1152--1162). The article's complete original proof of it is delivered with it (source0.tex lines 1178--1259, one \texttt{proof} environment of 82 lines, which the article places away from the statement). No mathematics is written, repaired or re-derived: the statement and the proof are the article's own lines, in the article's own notation, and no retained line is substituted or re-flowed.  Retained uncounted as the source-local context the proof needs: lines 286--288; lines 309--310; lines 313--323; lines 635--646; lines 754--759; lines 760--761; lines 763--765; lines 787--792; lines 793--798; lines 918--929; lines 1164--1169.  Nothing was omitted from any retained span, so the package carries no omission record.  No retained line contains a citation, so the wrapper carries no bibliography and no bibliography entry.  No retained line contains a figure environment, an image-inclusion command or drawing code, so no image is delivered and the package declares no asset dependency.  Editorial formatting only, in addition: an AMS \texttt{amsart} preamble that restates the article's own declarations and macros used by the retained lines, namely the environments \texttt{assumption}, \texttt{proposition}, \texttt{remark}, \texttt{lemma}, \texttt{theorem}, \texttt{corollary} on separate counters, together with the standard packages those lines need.  The retained environments print in this document as Assumption 1, Proposition 1, Remark 1, Lemma 1, Theorem 1, Corollary 1 in order of appearance, whereas the article prints the same items with the numbering of its own full document.  The complete native source of the article as distributed in the arXiv e-print for this exact version is bundled unchanged as \texttt{sources/194073/source0.tex}.
\begin{assumption}\label{assump:data} Data points $\{x_0^{(i)}\}_{i=1}^n$ are well-separated and bounded, i.e., $\Delta\coloneqq \min_{i\neq j} \| x_0^{(i)}-x_0^{(j)} \|>0$ and $R \coloneqq \max_i \| x_0^{(i)} \|<\infty$.
The data distribution $p$ is the \emph{empirical} distribution of the data, i.e., $p= \tfrac{1}{n}\sum_{i=1}^n \delta_{x_0^{(i)}}$. 
\end{assumption}

\begin{assumption}[Smoothness and positive reach]\label{assump:ridge set} For any $t\in [\delta, T]$, there exists $r_t>0$ such that $\mc{R}_t$ is a (piecewise) $C^2$-embedded submanifold in $\mb{R}^d$ with a reach no smaller than $r_t>0$, i.e., for all $x\in \mb{R}^d$ with $\mathrm{dist}(x,\mc{R}_t)\le r_t$, there exists a unique nearest point on $\mc{R}_t$.
\end{assumption}

\begin{proposition}\label{prop:projection-regularity} Under Assumption \ref{assump:data}, the log-density ridge family $\{\mc{R}_t\}_{\delta\le t\le T}$ satisfies Assumption \ref{assump:ridge set}. More precisely, as $t\to \delta^+\ll 1$, the reach satisfies $r_t = \Omega(h_t^{2} \theta_t^{-1} R^{-3})$ for arbitrary $\theta_t=\exp(- o(h_{t}^{-1}) )$. For any radius $\rho_t\in (0,r_t)$, define the tube neighborhood
   \begin{align}\label{eq:tube ridge}
    \mc{T}_t(\rho_t)\coloneqq \{ x\in \mb{R}^d | \mathrm{dist}(x, \mc{R}_t)\le \rho_t \}.
\end{align} 
Then the nearest-point projection $\Pi_t: \mc{T}_t(\rho_t)\to \mc{R}_t$ is well-defined, and
\begin{itemize}
    \item [(1)] for all $ x\in \mc{T}_t(\rho_t)$, the displacement $ x-\Pi_t(x)$ lies in the normal space of $\mc{R}_t$ at $\Pi_t(x)$;
    \item [(2)]  $\Pi_t$ is $C^1$ on $\mc{T}_t(\rho_t)$ and $\sup_{x\in \mc{T}_t(\rho_t)} \| \nabla \Pi_t(x) \|\le \tfrac{1}{1-\rho_t/r_t}$;
    \item [(3)] if $\rho_t=\Theta(r_t)$, the ridge motion is uniformly bounded: $\sup_{x\in \mc{T}_t(\rho_t)}\| \partial_t\Pi_t(x)\| = \mc{O}\left( R\right)$.
\end{itemize}
\end{proposition}

\begin{proposition}\label{prop:tube well-defined} Under Assumption \ref{assump:ridge set}, for any $\rho_t\in (0,r_t]$ and the tube neighborhood $\mc{T}_t(\rho_t)$ given below
   \begin{align}\label{eq:tube ridge 2}
    \mc{T}_t(\rho_t)\coloneqq \{ x\in \mb{R}^d | \mathrm{dist}(x, \mc{R}_t)\le \rho_t \},
\end{align} 
the nearest-point projection $\Pi_t: \mc{T}_t(\rho_t)\to \mc{R}_t$ is well-defined and we have
\begin{itemize}
    \item [(1)] $\Pi_t$ is $C^1$ on $\mc{T}_t(\rho_t)$;
    \item [(2)] $\forall x\in \mc{T}_t(\rho_t)$, $n_t(x)\coloneqq x-\Pi_t(x)$ is in the normal space at $\Pi_t(x)$, i.e, $n_t(x) \in N_{\Pi_t(x)}(\mc{R}_t)$;
    \item [(3)]  $\sup_{x\in \mc{T}_t(\rho_t)} \| \nabla \Pi_t(x) \|\le \tfrac{1}{1-\rho_t/r_t}$;
    \item [(4)] the motion of $\Pi_t$ is bounded: for any $ z\in \mc{R}_t$, there exists a velocity field $v_t\in N_z(\mc{R}_t)$ s.t. $\sup_{x\in \mc{T}_t(\rho_t)} \| \partial_t \Pi_t\| \le V_t + \frac{\rho_t}{1-\rho_t/r_t} W_t$ where $$V_t=\sup_{z\in \mc{R}_t} \| v_t(z) \|,\quad W_t= \sup_{z\in \mc{R}_t, \| u \|=1} \| P^\parallel(z)(\nabla v_t(z)u) \|\footnote{For any $z\in \mc{R}_t$, we use $P^\parallel(z)$ (or $P^\perp(z)$) to represent the orthogonal projection from $\mb{R}^d$ to $T_z(\mc{R}_t)$ (or $N_z(\mc{R}_t)$).}. $$
\end{itemize}
\end{proposition}

\begin{proposition}\label{prop:constant estimation} Under Assumption \ref{assump:data}, the log-density ridge sets satisfy Assumption \ref{assump:ridge set}. Furthermore, when $t\to \delta^+\ll 1$, we have the following order estimations:
\begin{align*}
    r_t = \Omega(\beta_t h_t^{3} \theta_t^{-1} R^{-3}),\quad V_t = \mc{O}( {\beta_t^{-1} h_t^{-1}} R  ),\quad W_t =\mc{O} ( \beta_t^{-1}h_t^{-1} R r_t^{-1} ),
\end{align*}
where $\theta_t$ is arbitrarily with order $\theta_t=\exp(- o(h_{t}^{-1}) )$.
\end{proposition}

\begin{remark}[Order estimation of ridge motion]\label{rem:ridge motion bound} Combining Propositions \ref{prop:tube well-defined} and \ref{prop:constant estimation}, picking $\rho_t=\Theta(r_t)$ and $\beta_t=\Theta({1}/{h_t})$, we proved Proposition \ref{prop:projection-regularity}-(3): $\sup_{x\in \mc{T}_t(\rho_t)}\| \partial_t\Pi_t(x)\| = \mc{O}\left( R\right)$. 
\end{remark}

\begin{align}\label{eq:data dominate region}
    \mc{B}_{T-t}^{(i)}(\theta_{T-t}) \coloneqq \{ x\in \mb{R}^d \mid \Softmax(-\frac{\| x-a_{T-t}x_0 \|^2}{2h_{T-t}})_i \ge 1-\theta_{T-t} \}.
\end{align}

\begin{lemma}\label{lem:important estimate} For any $t\in [\delta,T]$, we have
\begin{align*}
    \sup_{x\in \cup_{i=1}^n \mc{B}_{t}^{(i)}(\theta_{t})} \| \nabla m(t,x) \|\le \frac{20a_t^2 \theta_t R^2}{h_t},\quad  \sup_{x\in \cup_{i=1}^n \mc{B}_{t}^{(i)}(\theta_{t})}\| \nabla^3 \log p_t(x) \|\le \frac{80 a_t^3\theta_t R^3}{h_t^{3}},
\end{align*} 
where $\mc{B}_{t}^{(i)}(\theta_{t})$ is defined in \eqref{eq:data dominate region} with any $\theta_t=\exp(- o(h_{t}^{-1}) )$.
\end{lemma}

\begin{remark}[Choice of $\beta_t$]\label{rem:choice of beta} Lemma \ref{lem:important estimate} also validates the choice of ridge threshold $\beta_t=\Theta(1/h_t)$ when $t\to 0^+$. According to Lemma \ref{lem:explicit formula}, $\nabla^2\log p_t(x) = -\tfrac{1}{h_t}I_d + \tfrac{1}{h_t^2}\Sigma(t,x)$. According to estimations in Lemma \ref{lem:important estimate}, each eigenvalue $\lambda_j(t,x)$ of $\nabla^2\log p_t(x)$ satisfies
\begin{align*}
    \lambda_j(t,x) \le -\frac{1}{h_t} + \frac{20a_t^2 \theta_t R^2 }{h_t^2} \le -\frac{c}{h_t}, \quad \frac{1}{2}<c<1,
\end{align*}
under some choice of $\theta_t=\exp(- o(h_{t}^{-1}) )$. Therefore, as $t\to 0^+$, the choice of $\beta_t=c/h_t$ makes the second condition in the log-density ridge definition automatically satisfied.    
\end{remark}

\begin{lemma}[Explicit formulas]\label{lem:explicit formula} Under Assumption \ref{assump:data}, we have that for all $t\in [\delta,T]$,
    \begin{align*}
        \nabla \log p_t(x) &= -\frac{x}{h_t}+\frac{m(t,x)}{h_t}, \quad \nabla^2 \log p_t(x) = -\frac{I_d}{h_t}+\frac{\Sigma(t,x)}{h_t^2},\\
        \nabla^3\log p_t(x) &= \frac{\mb{E}[(U(t,x)-m(t,x))^{\otimes 3}]}{h_t^3},
    \end{align*}
    where $U(t,x)\in \mb{R}^d$ is a random vector taking values $\{a_tx_0^{(i)}\}_{i=1}^n$ with probabilities $\{\Softmax( -\frac{\| x-a_t x_0 \|^2}{2h_t} )_i\}_{i=1}^n$ and 
    \begin{align*}
        m(t,x) &= \mb{E}[U(t,x)], \quad
        \Sigma(t,x)  =\mathrm{Cov}(U(t,x)),
    \end{align*}
    with $x_0=(x_0^{(1)}, x_0^{(2)},\cdots, x_0^{(n)})$. Furthermore, $\nabla m(t,x)= \Sigma(t,x)/h_t$.
\end{lemma}

\begin{corollary}\label{cor:distance} As a consequence of Proposition \ref{prop:tube well-defined}, the following properties hold for the square-distance function $D_t(x)\coloneqq \| x-\Pi_t(x) \|^2$:
    \begin{itemize}
        \item [(1)] $\mc{D}_t\in C^2(\mc{T}_t(\rho_t))$ and $\nabla D_t(x)=2n_t(x)$;
        \item [(2)] $\sup_x \| \nabla^2 D_t(x) \|\le \tfrac{2}{1-\rho_t/r_t}$. As a consequence, $\sup_x \Delta D_t(x)\le \tfrac{2d}{1-\rho_t/r_t}$.
    \end{itemize}
\end{corollary}

\begin{theorem}\label{thm: formal normal contraciton} Under Assumption \ref{assump:data}, define $\kappa_{s,t}=2\int_s^t (\beta_{T-u}-1)\dee u$, $B_{s,t}(d,R)=\int_{s}^t e^{-\kappa_{u,t} } ( \rho_{T-u}R + d ) \dee u$ and $e_A^\perp(t,x)\coloneqq P^\perp(\Pi_t(x))e_A(t,x)$ with $e_A=m_A-m$. Then for any $t\in [\Tilde{t}_{\mathrm{in}},T-\delta]$,
{
    \begin{align*}
    \mb{E}[D_{T-t}(\Tilde{Y}_t)] &\lesssim e^{-\kappa_{\Tilde{t}_{\mathrm{in}},t} }\rho_{T-\Tilde{t}_{\mathrm{in}}}  + B_{\Tilde{t}_{\mathrm{in}},t}(d,R)  + \int_{\Tilde{t}_{\mathrm{in}}}^t e^{-\kappa_{u,t} } \frac{\mb{E}[\| e_A^\perp(T-u,\Tilde{Y}_u)\|^2]}{h_{T-u}^2\beta_{T-u}} \dee u. 
    \end{align*}}
    Furthermore, picking picking $\beta_t=c/h_t$ for any $c\in [\frac{1}{2},1)$\footnote{$c$ can be chosen arbitrarily between $[\tfrac{1}{2},1)$ due to the property of $\nabla^2\log p_t(x)$ as explained in Remark~\ref{rem:choice of beta}. }, when $\delta\ll 1$, we have
    {
\begin{align*}
    \mb{E}[D_{\delta}(\Tilde{Y}_{T-\delta})] = \mc{O}(d\delta^c+ \delta^c\int_{\Tilde{t}_{\mathrm{in}}}^{T-\delta}  h_{T-u}^{-1-c}\,{\mb{E}[\| e_A^\perp(T-u,\Tilde{Y}_u)\|^2]} \dee u  ).
\end{align*}}
\end{theorem}

\begin{proof}[Proof of Theorem \ref{thm: formal normal contraciton}]
For any $t\in [0,T-\delta]$, $x\in \mc{T}_{T-t}(\rho_t)$, recall that $D_{T-t}(x)=\| n_{T-t}(x) \|^2$ and $n_{T-t}(x)=x-\Pi_{T-t}(x)$. Once the processes enter $\mc{T}_{T-t}(\rho_t)$, we can track the evolution of $D_{T-t}$ via It\^o's formula. For the reverse process $Y_t$, we have
\begin{align*}
    \dee D_{T-t}(Y_t) &= \big( \partial_t D_{T-t}(Y_t) + 2\langle n_{T-t}(Y_t), b(T-t,Y_t) \rangle + 2\Tr(\nabla^2 D_{T-t}(Y_t))\big)\dee t \\
    &\quad + 2\sqrt{2}\langle n_{T-t}(Y_t),\dee \bar{B}_t \rangle,
\end{align*}
where we have the following estimations for the drift terms:
\begin{itemize}
    \item [(1)] According to Corollary \ref{cor:distance}, $2\Tr(\nabla^2 D_{T-t}(Y_t))\le \tfrac{4d}{1-\rho_{T-t}/r_{T-t}}$. Picking $\rho_t=r_t/2$ for all $t$ and we get $2\Tr(\nabla^2 D_{T-t}(Y_t))\le 8d$.
    \item [(2)] According to Remark \ref{rem:ridge motion bound}, 
    \begin{align*}
        \partial_t D_{T-t}(Y_t)&=-\partial_s D_s(Y_t)|_{s=T-t}= 2\langle n_{T-t}(Y_t), \partial_s \Pi_s (Y_t) |_{s=T-t}  \rangle \\
        &\le 2\rho_{T-t} \sup_{x\in \mc{T}_{T-t}(\rho_{T-t})}\| \partial_s\Pi_s(x)|_{s=T-t} \|  \coloneqq 2\rho_{T-t} S_{T-t} ,
    \end{align*}
    and $S_{T-t}= \mc{O}(R)$ when $T-t\to 0^+$.
    \item [(3)] if we denote $z=z_{T-t}\coloneqq \Pi_{T-t}(Y_t)$, $n_{T-t}(Y_t)\perp T_z(\mc{R}_{T-t})$ and we have
    \begin{align*}
         2\langle n_{T-t}(Y_t) , b(T-t,Y_t) \rangle  & = 2\langle  n_{T-t}(Y_t) ,  Y_t \rangle + 4 \langle  n_{T-t}(Y_t) , \nabla\log p_{T-t}(Y_t) \rangle \\
        & = 2\| n_{T-t}(Y_t) \|^2 + 4 \langle  n_{T-t}(Y_t) , \nabla\log p_{T-t}(Y_t) \rangle,
    \end{align*}
    In the last term, we can write $Y_t=z+n_{T-t}(Y_t)$ and expand $\nabla\log p_{T-t}(Y_t)$ at $z\in \mc{R}_{T-t}$. Then we have
    \begin{align*}
        &\quad \langle  n_{T-t}(Y_t) , \nabla\log p_{T-t}(Y_t) \rangle \\
        & = \langle  n_{T-t}(Y_t) ,  \nabla\log p_{T-t}(z) \rangle + \langle  n_{T-t}(Y_t) ,  \nabla^2\log p_{T-t}(z) n_{T-t}(Y_t) \rangle \\
        & \quad  + \frac{1}{2}\langle  n_{T-t}(Y_t) ,  \nabla^3\log p_{T-t}(z') n_{T-t}(Y_t)^{\otimes 2} \rangle ,
    \end{align*}
    where
    \begin{align*}
       & \langle  n_{T-t}(Y_t) ,  \nabla\log p_{T-t}(z) \rangle = 0 ,\quad &\text{definition of }\mc{R}_{T-t},\\
       & \langle  n_{T-t}(Y_t) ,  \nabla^2\log p_{T-t}(z) n_{T-t}(Y_t) \rangle  \quad & \\
       = & \langle  n_{T-t}(Y_t) ,  P^\perp(z)\nabla^2\log p_{T-t}(z) n_{T-t}(Y_t) \rangle \le  -\beta_{T-t}\| n_{T-t}(Y_t) \|^2, \quad & \text{definition of }\mc{R}_{T-t},\\ 
       &\langle  n_{T-t}(Y_t) ,  \nabla^3\log p_{T-t}(z') n_{T-t}(Y_t)^{\otimes 2} \rangle \quad & \\
       \le & \sup_x \| \nabla \log p_{T-t}(x) \| \| n_{T-t}(Y_t) \|^3 \le  \frac{80a_{T-t}^3 R^3}{h_{T-t}^3} \rho_{T-t} \| n_{T-t}(Y_t) \|^2 \quad & \text{Lemma }\ref{lem:important estimate} \\
       \le & \frac{a_{T-t}^3\beta_{T-t}}{2} D_{T-t}(Y_t) \quad & \text{Proposition }\ref{prop:constant estimation}
    \end{align*}
    The last identity follows from Proposition \ref{prop:constant estimation} by picking $\theta_t$ such that $r_t= \beta_t h_t^3R^{-3}/80$ and $\rho_t=r_t/2$. Combining the above inequalities, we have
    \begin{align*}
        \langle  n_{T-t}(Y_t) , \nabla\log p_{T-t}(Y_t) \rangle &\le -\beta_{T-t}(1-a_{T-t}^3/4) \| n_{T-t}(Y_t) \|^2\le -\frac{3}{4}\beta_{T-t} D_{T-t}(Y_t).
    \end{align*}
    Therefore, we have
\begin{align*}
    2\langle n_{T-t}(Y_t), b(T-t,Y_t) \rangle  \le -(3\beta_{T-t}-2)D_{T-t}(Y_t).
\end{align*}
\end{itemize}
Combining all the estimations and taking expectations of $D_{T-t}$, we obtain the following inequality
\begin{align*}
    \frac{\dee}{\dee t}\mb{E}[D_{T-t}(Y_t)] \le -(3\beta_{T-t}-2)\mb{E}[D_{T-t}(Y_t)] + 2\rho_{T-t}S_{T-t} + 8d.
\end{align*}
Last, apply Gronwall's inequality, for any $\tin\in (0, T-\delta)$ and $t\in (\tin,T-\delta]$, we obtain
\begin{align*}
    \mb{E}[D_{T-t}(Y_t)] &\le \exp\big( -\int_{\tin}^t 3\beta_{T-u}-2 \dee u \big)\mb{E}[D_{t-\tin}(Y_{\tin})]  \\
    &\quad + \int_{\tin}^t \exp\big( -\int_{u}^t 3\beta_{T-s}-2 \dee s \big) \big( 2\rho_{T-u}S_{T-u} + 8d \big) \dee u .
\end{align*}
When $t=T-\delta$, $\beta_t=c/h_t$ and $\delta\ll 1$, we have
\begin{align*}
    \int_{\tin}^{T-\delta} \exp\big( -\int_{u}^{T-\delta
    } 3\beta_{T-s}-2 \dee s \big) \big( 2\rho_{T-u}S_{T-u} + 8d \big) \dee u= \mc{O}(d\delta^{\frac{3c}{2}}).
\end{align*}   

\noindent For the reverse-time inference process $\Tilde{Y}_t$, we have
\begin{align*}
    \dee D_{T-t}(\Tilde{Y}_t) &= \big( \partial_t D_{T-t}(\Tilde{Y}_t) + 2\langle n_{T-t}(\Tilde{Y}_t), b_A(T-t,\Tilde{Y}_t) \rangle \\
    &\quad + 2\Tr(\nabla^2 D_{T-t}(\Tilde{Y}_t))\big)\dee t + 2\sqrt{2}\langle n_{T-t}(\Tilde{Y}_t),\dee \Tilde{B}_t \rangle.
\end{align*}
Notice that the only difference to that of $Y_t$ is the error $e_A(t,\cdot)=m(t,\cdot)-m_A(t,\cdot)$ within the normal space, i.e., $e_A^\perp(t,x)\coloneqq P^\perp(\Pi_t(x)) e_A(t,x)$. Similarly, we have the inequality
\begin{align*}
     &\quad \frac{\dee}{\dee t}\mb{E}[D_{T-t}(\Tilde{Y}_t)] \\
     &\le -(3\beta_{T-t}-2)\mb{E}[D_{T-t}(\Tilde{Y}_t)] + 2\rho_{T-t}S_{T-t} + 8d + 2\mb{E}[\langle n_{T-t}(\Tilde{Y}_t), \frac{1}{h_{T-t}}e_A(T-t,\Tilde{Y}_t) \rangle]\\
      &\le -2(\beta_{T-t}-1)\mb{E}[D_{T-t}(\Tilde{Y}_t)] + 2\rho_{T-t}S_{T-t} + 8d + \frac{\mb{E}[\| e_A^\perp(T-t,\Tilde{Y}_t)\|^2]}{h_{T-t}^2\beta_{T-t}},
\end{align*}
where the last inequality follows from Young's inequality. Therefore, according to the Gronwall's inequality,
\begin{align*}
    \mb{E}[D_{T-t}(\Tilde{Y}_t)] & \le  \exp\big( -2\int_{\Tilde{t}_{\mathrm{in}}}^t \beta_{T-u}-1 \dee u \big)\mb{E}[D_{t-\Tilde{t}_{\mathrm{in}}}(\Tilde{Y}_{\Tilde{t}_{\mathrm{in}}})] \\
    &\quad + \int_{\Tilde{t}_{\mathrm{in}}}^t \exp\big( -2\int_{u}^t \beta_{T-s}-1 \dee s \big) \big( 2\rho_{T-u}S_{T-u}+\frac{\mb{E}[\| e_A^\perp(T-u,\Tilde{Y}_u)\|^2]}{h_{T-u}^2\beta_{T-u}}+ 8d \big) \dee u .
\end{align*}
When $t=T-\delta$, $\beta_t=c/h_t$ and $\delta\ll 1$, we have
\begin{align*}
    &\quad\int_{\Tilde{t}_{\mathrm{in}}}^{T-\delta} \exp\big( -2\int_{u}^{T-\delta} \beta_{T-s}-1 \dee s \big) \dee u =\mc{O}(\delta^c) ,\\
    &\quad\int_{\Tilde{t}_{\mathrm{in}}}^{T-\delta} \exp\big( -2\int_{u}^{T-\delta} \beta_{T-s}-1 \dee s \big) h_{T-u}^{-2}\beta_{T-u}^{-1}\mb{E}[\| e_A^\perp(T-u,\Tilde{Y}_u)\|^2] \dee u \\
    & =\mc{O}(\delta^c \int_{\Tilde{t}_{\mathrm{in}}}^{T-\delta}  h_{T-u}^{-1-c}\mb{E}[\| e_A^\perp(T-u,\Tilde{Y}_u)\|^2] \dee u ) .
\end{align*}
\end{proof}

\end{document}
